\documentclass[12pt]{article}
\usepackage[utf8]{inputenc}
\usepackage[T1]{fontenc}
\usepackage{hyperref}
\title{Yuqui Language Resources for K-12 Curriculum Development}
\author{Compiled by Claude Sonnet 5.5 (claude-sonnet-5-5) for Andr\'as Kornai}
\date{}
\begin{document}
\maketitle

\section{Introduction}
Yuqui (Mbiaye', ISO 639-3 yuq, Glottocode yuqu1240) is a severely endangered Tupi-Guaran\'i language of Cochabamba department, Bolivia. Speaker estimates in the sources found range from about 120 to about 140. This report rests on a single search. Very little was found, and no school curriculum or classroom materials for Yuqui turned up.

\section{Primary Educational Resources}
No primers, textbooks or school programmes for Yuqui were found in the search results.

\section{Linguistic and Reference Materials}
\begin{itemize}
\item Villafa\~ne (2004), \emph{Gram\'atica Yuki lengua Tupi-Guaran\'i de Bolivia}, a Spanish-language thesis from Nijmegen. It is the most substantial resource found. A PDF is listed among the references of the Wikipedia article: \url{https://en.wikipedia.org/wiki/Yuqui_language}
\item Omniglot page with an alphabet chart and pronunciation notes: \url{https://omniglot.com/writing/yuqui.htm}
\item Glottolog entry: \url{https://glottolog.org/resource/languoid/id/yuqu1240}
\end{itemize}
No dictionary was found in the results.

\section{Religious and Community Resources}
A partial Bible translation into Yuqui (2000) is mentioned on the Omniglot page above. No copy was located in the search results.

\section{Digital and Audio Resources}
The Endangered Languages Project has entries for the language: \url{https://endangeredlanguages.com/node/37551}

\section{Curriculum Development Context}
The documentation found is thin: one grammar and a few reference pages. A curriculum would need new work with the community. Any further resources, such as the SIL materials mentioned on Omniglot, should be checked directly. They were not verified here.
\end{document}